\section{Konsep Dasar Block Cipher}

\textit{Block cipher} atau \textbf{cipher blok} adalah primitif simetris yang
bekerja dengan membagi plainteks menjadi blok-blok bit yang panjangnya sama,
lalu mengenkripsi setiap blok sebagai satu kesatuan. Bila pada
Bab~\ref{chap:stream-cipher} cipher alir memproses satu bit atau satu byte
setiap kali, cipher blok justru menunggu terkumpulnya satu blok penuh
--- umumnya 64 bit atau 128 bit --- sebelum operasi enkripsi dijalankan.
Perbedaan satuan kerja inilah yang membedakan kedua keluarga cipher simetris
dan yang menentukan di mana masing-masing dipakai.

Kedua pendekatan itu menyelesaikan persoalan yang berbeda. Cipher alir
unggul ketika data tiba terus-menerus dan harus dienkripsi tanpa penundaan,
misalnya pada percakapan suara. Cipher blok unggul ketika data tersedia
utuh dan harus dilindungi dalam satuan yang dapat diakses sendiri-sendiri,
misalnya berkas, rekaman basis data, atau blok data di dalam protokol
jaringan. Karena sifatnya yang dapat diproses per blok, cipher blok juga
menjadi bahan dasar bagi konstruksi lain --- fungsi \textit{hash},
kode otentikasi pesan, dan mode operasi terautentikasi --- sebagaimana akan
terlihat pada Bab~\ref{chap:hash-mac}.

\subsection{Karakteristik Utama}

Tiga ciri berikut membedakan cipher blok dari cipher alir.

\begin{enumerate}
    \item \textbf{Ukuran blok tetap.} Plainteks dibagi menjadi blok berukuran
          $n$ bit, dengan $n$ yang lazim dipakai adalah 64 bit, 128 bit, atau
          256 bit. Ukuran blok tidak berubah selama enkripsi maupun dekripsi
          berlangsung.
    \item \textbf{Panjang keluaran sama dengan panjang masukan.} Panjang blok
          cipherteks selalu sama persis dengan panjang blok plainteksnya,
          yaitu $n$ bit. Tidak ada penambahan bit di dalam blok, hanya
          penambahan \textit{padding} pada blok terakhir bila pesan tidak
          habis dibagi ukuran blok.
    \item \textbf{Kunci eksternal bebas panjangnya.} Panjang kunci yang
          diberikan pengguna tidak harus sama dengan panjang blok. Pada DES
          panjang blok 64 bit tetapi kunci efektifnya 56 bit; pada AES-256
          panjang blok 128 bit sedangkan kuncinya 256 bit.
\end{enumerate}

Secara matematis, bila $P$ adalah sebuah blok plainteks berukuran $n$ bit
dan $C$ adalah blok cipherteks hasilnya, maka
\begin{equation}
P = (p_1, p_2, \dots, p_n), \quad p_i \in \{0,1\},
\label{eq:blok-plainteks}
\end{equation}
\begin{equation}
C = (c_1, c_2, \dots, c_n), \quad c_i \in \{0,1\},
\label{eq:blok-cipherteks}
\end{equation}
dan proses enkripsi serta dekripsinya dituliskan sebagai
\begin{equation}
C = E_K(P) \qquad \text{dan} \qquad P = D_K(C),
\label{eq:block-encrypt-decrypt}
\end{equation}
dengan $E$ fungsi enkripsi, $D$ fungsi dekripsi, dan $K$ kunci. Karena
keduanya bekerja pada blok dengan panjang yang sama, $E_K$ dan $D_K$ sama-sama
memetakan himpunan $n$ bit ke himpunan $n$ bit.

\subsection{Block Cipher sebagai Permutasi}

Ada satu syarat yang tidak boleh dilanggar oleh setiap cipher blok, dan
syarat itu membatasi seperti apa algoritmanya boleh dibentuk. Untuk setiap
kunci $K$, fungsi $E_K$ harus \textbf{bijektif}: setiap blok cipherteks
$C$ harus berasal dari tepat satu blok plainteks $P$, yaitu $P = D_K(C)$.
Bila dua plainteks berbeda $P_1 \neq P_2$ dapat menghasilkan cipherteks
yang sama, yaitu $E_K(P_1) = E_K(P_2)$, maka penerima tidak punya cara
untuk memutuskan plainteks mana yang sebenarnya dikirim --- dekripsi menjadi
tidak mungkin dilakukan.

Akibatnya, untuk setiap kunci tetap, $E_K$ merupakan sebuah \textbf{permutasi}
dari himpunan $2^n$ blok yang mungkin. Pemetaan inilah yang disebut
\textbf{permutasi semu} (\textit{pseudo-random permutation}) karena sifatnya
yang harus tampak acak bagi penyerang, meskipun sesungguhnya ditentukan
sepenuhnya oleh kunci. Himpunan semua permutasi atas $n$ bit berjumlah
$2^n!$ buah, jauh lebih besar daripada jumlah kunci yang mungkin, yaitu
$2^m$ untuk kunci sepanjang $m$ bit. Jadi sebuah cipher blok hanya memilih
$2^m$ permutasi dari $2^n!$ permutasi yang tersedia, dan pilihan itu
dikendalikan oleh kunci.

Pandangan ini menjelaskan dua hal sekaligus. Pertama, ia menunjukkan mengapa
cipher blok dapat dibalik tanpa menyimpan tabel apa pun: yang perlu diketahui
penerima hanyalah kunci yang memilih permutasi tersebut. Kedua, ia menjelaskan
mengapa kunci harus panjang: dengan $m$ bit kunci, penyerang yang menebak
kunci secara menyeluruh hanya perlu mencoba $2^m$ kemungkinan, berapa pun
besar $n$. Untuk blok 128 bit dengan kunci 128 bit, ruang kunci berukuran
$2^{128} \approx 3{,}4 \times 10^{38}$, dan jumlah itu jauh melampaui
kemampuan komputasi yang ada saat ini.

\subsection{Ukuran Blok, Kunci, dan Padding}

Ukuran blok dan ukuran kunci berperan berbeda terhadap keamanan. Ukuran blok
yang besar memperluas ruang yang harus dipetakan oleh permutasi, sehingga
\textit{diffusion} --- penyebaran pengaruh satu bit ke banyak bit lain ---
menjadi lebih efektif dan pola statistik lebih sulit muncul. Ukuran kunci
yang besar menambah jumlah permutasi yang dapat dipilih, sehingga pencarian
kunci menyeluruh menjadi lebih mahal. Tabel~\ref{tab:ukuran-blok-kunci}
meringkas parameter beberapa cipher blok yang dibahas di dalam buku ini.

\begin{table}[htbp]
\centering
\small
\begin{tabular}{@{}lcc c@{}}
\hline
\textbf{Algoritma} & \textbf{Ukuran blok} & \textbf{Ukuran kunci} & \textbf{Putaran} \\
\hline
DES            & 64 bit  & 56 bit            & 16 \\
DES tiga kunci & 64 bit  & 112 atau 168 bit  & 48 \\
AES-128        & 128 bit & 128 bit           & 10 \\
AES-192        & 128 bit & 192 bit           & 12 \\
AES-256        & 128 bit & 256 bit           & 14 \\
GOST           & 64 bit  & 256 bit           & 32 \\
Blowfish       & 64 bit  & 32--448 bit       & 16 \\
\hline
\end{tabular}
\caption{Ukuran blok, ukuran kunci, dan jumlah putaran beberapa cipher blok}
\label{tab:ukuran-blok-kunci}
\sumber{Data dari Munir, \textit{IF4020 Kriptografi}, ITB; dilengkapi dari \citet{stallings2017}.}
\end{table}

Bila panjang pesan tidak habis dibagi oleh ukuran blok, blok terakhir menjadi
lebih pendek daripada blok-blok lainnya. Blok yang tidak lengkap itu harus
ditambahi bit \textbf{\textit{padding}} agar panjangnya kembali sama dengan
$n$ bit. Bit padding dapat berupa bit nol semua, bit satu semua, atau bit nol
dan satu yang berselang-seling. Sebagai contoh, dengan ukuran blok 8 bit,
pesan berikut
\begin{center}
\texttt{10001101\ 00101001\ 10110}
\end{center}
memiliki blok terakhir yang hanya 5 bit, sehingga perlu ditambahi tiga bit
padding, misalnya:
\begin{center}
\texttt{10001101\ 00101001\ 10110}\textbf{\texttt{000}}.
\end{center}
Cara penerima membedakan bit padding dari bit pesan yang sebenarnya ---
dan karena itu cara membuang padding ketika dekripsi --- dibahas bersama
mode operasi pada Bagian~\ref{sec:mode-operasi}.

\subsection{Contoh Sederhana Block Cipher}

Untuk memperlihatkan bagaimana enkripsi dan dekripsi bekerja pada satuan blok,
tinjau sebuah cipher blok mainan berukuran 4 bit. Fungsi enkripsinya melakukan
XOR antara plainteks $P$ dengan kunci $K$, lalu menggeser bit-bit hasilnya satu
posisi ke kiri secara siklik (\textit{wrapping}):
\begin{equation}
E_K(P) = (P \oplus K) \ll 1.
\label{eq:mainan-enkripsi}
\end{equation}
Karena dekripsi harus membalik urutan operasi itu, dekripsinya adalah menggeser
siklik cipherteks satu posisi ke kanan terlebih dahulu, baru di-XOR dengan kunci:
\begin{equation}
D_K(C) = (C \gg 1) \oplus K.
\label{eq:mainan-dekripsi}
\end{equation}
Kedua rumus itu memang saling membalik. Menggeser $C = (P \oplus K) \ll 1$
satu posisi ke kanan mengembalikan $P \oplus K$ dengan utuh, dan meng-XOR hasil
itu dengan $K$ memakai sifat $a \oplus a = 0$ pada
Bab~\ref{chap:stream-cipher} sehingga menyisakan $P$.

Sebagai contoh, ambil $P = \code{1010}$ dan $K = \code{1011}$. Maka
$P \oplus K = \code{0001}$, dan geser siklik ke kiri menghasilkan
$C = \code{0010}$. Sisi penerima menggeser $\code{0010}$ ke kanan menjadi
$\code{0001}$, lalu meng-XOR-nya dengan $K = \code{1011}$; hasilnya kembali
$\code{1010}$, yaitu plainteks semula. Contoh mainan ini sengaja dipilih sesederhana
mungkin, tetapi ia sudah memperlihatkan dua hal yang berlaku umum: enkripsi dan
dekripsi adalah proses yang saling membalik, dan penambahan satu blok data dapat
mengubah satu blok penuh cipherteks.

Contoh ini juga sudah cukup untuk menunjukkan mengapa cipher blok tidak boleh
dipakai apa adanya. Karena $E_K$ bekerja pada blok yang sama secara
independen, dua blok plainteks yang identik akan selalu menghasilkan dua blok
cipherteks yang identik pula. Kelemahan itu tidak berasal dari fungsi $E_K$
yang lemah, melainkan dari cara blok-blok dirangkai, dan cara
mengatasinya adalah pokok bahasan bagian berikutnya.
